\subsection{外延公理}

外延公理如下：

A1.

\[
(\forall x) (\forall y) ((x \subset y \text {   且   } y \subset x) \Longrightarrow (x = y)).
\]

直观上，这条公理表达了两个具有相同元素的集合是相等的这一事实。

为了证明 x = y，因此只需证明 \(z \in y\) 在通过附加该假设所得的理论中 \(z \in x\) ，以及 \(z \in x\) 在通过附加该假设所得的理论中 \(z \in y\) ，其中 z 是一个不同于 x、y 及常量的字母。

C48. 设 R 是一个关系，x 是一个字母，y 是一个不同于 x 且不出现在 R 中的字母。则关系 \((\forall x)((x \in y) \Longleftrightarrow R)\) 在 y 上是单值的。

让 \(\pmb{z}\) 是一个与 \(\pmb{x}\) 的字母，它不出现在 \(\pmb{R}\) 。添加假设

\[
(\forall x) ((x \in y) \Longleftrightarrow R), (\forall x) ((x \in z) \Longleftrightarrow R).
\]

于是我们依次有定理

\[
\begin{array}{c} (\forall x) (((x \in y) \Longleftrightarrow R) \quad \text { 且 } \quad ((x \in z) \Longleftrightarrow R)), \\ (\forall x) ((x \in y) \Longleftrightarrow (x \in z)), \\ y \subset z, \qquad z \subset y. \end{array}
\]

由 A1 我们得到 y = z。这证明了 C48。

